% Folder 134-2, batch 02 — archivists' pages 21–40.
% Pass: Opus 5.5 (claude-opus-5-5), 2026-10-03 — first pass, unchecked against the pages by a human.
%
% Special rules for folder 134-2 (issue #44): the Pursuing Stacks text (letter to
% Quillen and the numbered sections) is printed in Maltsiniotis (ed.), A la
% poursuite des champs, vol. I, SMF Documents mathematiques 20, 2022, and is not
% transcribed; Grothendieck's 1975 letters to Breen are transcribed in full.
%
% Contents of the batch:
% - pp. 21-36 (his pp. 7-22): letter to Daniel Quillen, 19.2-23.2.1983, English,
%   typescript annotated in his hand = Pursuing Stacks §§4 (end)-13 and PS of 23.2.
%   SKIPPED (published in Maltsiniotis vol. I). Dates on the pages: 21.2. (p. 28),
%   22.2.1983 (p. 34), 25.2 (marginal, p. 35), PS (23.2) (p. 36).
% - p. 37 (his p. 23): "APPENDIX to Chapter I: three letters to Larry Breen",
%   his English introduction (typescript, annotated). TRANSCRIBED. UNCLEAR whether
%   this belongs to Pursuing Stacks proper: it is in the PS typescript and
%   pagination, but it is neither the Quillen letter nor a numbered section, and
%   Maltsiniotis's preface defers the Breen material to a vol. II that has not
%   appeared; transcribed on that ground.
% - pp. 38-40 (his pp. 23 bis, 24, 25): first letter to Lawrence Breen, dated
%   "Villecun 5.2.1975", "Cher Breen", French, typed copy with his annotations
%   (paste-up: a retyped slip over the original typed copy, whose cancelled
%   opening shows through). TRANSCRIBED; runs on past p. 40 into batch 03.
% - No letter written by a third party in this batch.

\documentclass[11pt,a4paper]{article}
\input{../preamble/grothendieck}

\folder{134-2}
\batch{2}
\pages{21}{40}
\foldertitle{[Chapitre I : Take off (pages 1 à 65) et table des matières provisoire] : tapuscrits annotés (19/02-22/02/1983), note manuscrite (s.d.), copies de lettre (1975, s.d.).}
\dating{1975-[1983]}
\watermark{Édition de démonstration}

\begin{document}

\section{Lettre à Daniel Quillen (suite) — \emph{Pursuing Stacks}}

\page{21}
\note{Pp. 21 à 36 (pagination de l'auteur 7 à 22) : suite de la lettre à Daniel Quillen commencée le 19.2.1983, tapuscrit annoté de sa main ; c'est le texte de \emph{Pursuing Stacks}, §§ 4 (fin) à 13 et post-scriptum du 23.2 ; imprimé dans Maltsiniotis (éd.), À la poursuite des champs, vol. I, SMF, Documents mathématiques 20, 2022. Non transcrit. La fin du § 4 vient de la page précédente (lot 01).}

\note{p. de l'auteur 7 : fin du § 4, début du § 5.}

\page{22}
\note{p. de l'auteur 8 : § 5 (suite).}

\page{23}
\note{p. de l'auteur 9 : § 5 (suite).}

\page{24}
\note{p. de l'auteur 10 : fin du § 5, début du § 6.}

\page{25}
\note{p. de l'auteur 11 : § 6 (suite), début du § 7.}

\page{26}
\note{p. de l'auteur 12 : § 7 (suite).}

\page{27}
\note{p. de l'auteur 13 : § 7 (suite).}

\page{28}
\note{p. de l'auteur 14, datée « 21.2. » : §§ 8 et 9.}

\page{29}
\note{p. de l'auteur 15 : § 9 (suite).}

\page{30}
\note{p. de l'auteur 16 : fin du § 9, début du § 10.}

\page{31}
\note{p. de l'auteur 17 : § 10 (suite).}

\page{32}
\note{p. de l'auteur 18 : § 11.}

\page{33}
\note{p. de l'auteur 19 : § 12.}

\page{34}
\note{p. de l'auteur 20, datée « 22.2.1983 » : § 13.}

\page{35}
\note{p. de l'auteur 21 : § 13 (suite) ; une note marginale datée « 25.2 ».}

\page{36}
\note{p. de l'auteur 22 : fin du § 13, signature « Alexander », post-scriptum « PS (23.2) ».}

\section{Appendice au chapitre I : trois lettres à Larry Breen}

\page{37}
\note{p. de l'auteur 23. Texte en anglais, tapuscrit annoté de sa main. Il appartient à la dactylographie de \emph{Pursuing Stacks} mais n'est ni la lettre à Quillen ni une section numérotée ; voir l'en-tête du fichier.}

APPENDIX to Chapter I : three letters to Larry Breen

In this appendix, I am encluding\note{sic} three letters to Larry Breen, dated
5.2, 17.2 and 17-19.7. 1975. These letters are written in French, and the
first two are reproduced textually, whereas the third appears here in English
translation. I am grateful to Ronnie Brown who took the trouble last year to
make such a translation (from a hardly legible \struck{handwri} copy of the
handwritten letter to Larry Breen) with the assistance of Larry Breen himself
and J.L. Loday. I am now using his translation rather than the original
letter, which no printer could possibly decipher !

Also, for the second letter I am using a typed copy made in 1975 by Larry
Breen (who presumably had difficulties too deciphering the handwriting). Thanks
are due to him for his interest and patience (with someone like me, very
unknowledgeable in standard homotopy techniques), \struck{\ill{}} \add{which}
appeared \struck{from} \add{in} the letters I got from him in response and his
verbal explanations on related matters, \add{as well as \struck{from} in} the
trouble he took in \struck{copying} retyping \add{that letter} and sending me
copies of my own letters to him, and allowing me to reproduce them in this
volume of Pursuing Stacks. My only present contribution to this set of letters
is adding a few comments (in the notes (\add{22}) to (\add{31})),
\add{adding subtitles (with numbers App. 1 to App. 18), and correcting some}
\struck{and correcting some} inaccuracies in the English translation (due mainly
to my handwriting …). Also, I skipped the beginning of the first letter, which
does'nt seem \struck{to me} of general relevance.
\note{Après « (31)), » un court passage dactylographié est biffé sous l'ajout manuscrit ; illisible.}

The first two letters are an attempt to explain to Larry Breen (who has a wide
background in algebraic geometry and homological and homotopical algebra) some
of the main points of the programme I had in mind around the notions of
n-categories and n-stacks (which is what I am supposed to be pursuing now in my
present work "Pursuing Stacks"). They were written under the impetus of
\struck{the} \add{new} intuition \struck{which} (new to me at any rate) which
then had just appeared to me, namely that (non strict) n-groupoids should model
(in a suitable sense) n-truncated homotopy types.

\marginal{signe d'alinéa}
The third letter, \add{written in} answer to a number of questions in Larry
Breen's response to the first two, is of a wider scope. A large part of the
letter outlines (very sketchyly) some main points of a duality program
(including a cohomological formulation of "geometric" local and global
classfield theory), which \struck{appears} \add{emerged by the end of the
fifties and appears} here for the first time in print\struck{, and going back
to the end of the fifties}. \struck{It} \add{The later part of the letter}
gives also some hints about the need of a framework of "tame topology" suitable
for writing up a "dévissage theory" of stratified spaces, and for
\struck{\ill{}} \add{working} with étale tubular neighbourhoods, for the common
purpose of coming to grips \struck{of} \add{with} a suitable notion of "fine
homotopy type" of a "tame" topological space or \add{of} a scheme \add{say}, in
terms of the \struck{ordered set} \add{inverse system} of "indexed homotopy
types" corresponding to \add{all} equisingular stratifications.

\section{Première lettre à Lawrence Breen, Villecun, 5.2.1975}

\page{38}
\note{p. de l'auteur « 23 bis ». Copie dactylographiée, en français, annotée de sa main. Une bande retapée (date, appel « Cher Breen », coupure « …. » et les trois lignes de l'App. 1) est collée sur la copie dactylographiée d'origine ; le début de la lettre, sauté pour le volume (cf. p. 37), y transparaît sous la bande, barré de grandes croix. Ce qu'on en devine est donné ici, sous réserve ; le reste n'est pas lisible à travers le papier.}

\struck{Villecun le 5.2.1975 — Cher Breen, Merci pour ta longue lettre, reçue il y a
déjà un moment, et \uncertain{excuse} le retard pour répondre. Je dois avouer
que toute cette lettre m'a passé par dessus la tête, elle était écrite pour
quelqu'un \uncertain{plus} \ill{} que je ne le suis — surtout maintenant où je
ne réfléchis plus au maths que quelques jours par mois, à la sauvette.
\uncertain{Ton scepticisme} concernant la \ill{} exacte de ma lettre à Deligne
\uncertain{était fondé}, et ton interprétation \uncertain{cohomologique} \ill{}
comme un Ext (X(M),T) certainement exacte. Je m'étais convaincu par voie
géométrique qu'il y avait un homomorphisme d'obstruction \uncertain{Hom(M,N)
→ Ext (M,N)}, mais \ill{} au cas \uncertain{"universel"} \ill{}}

Villecun 5.2.1975

Cher Breen

….

(App. 1) … Pour tout le reste de ta lettre, elle mériterait un lecteur plus
averti, aussi, pour qu'elle ne soit pas entièrement perdue au monde, je vais
l'envoyer à Illusie ! J'ai néanmoins constaté, avec intérêt, ton intérêt à demi
refoulé pour des 2-catégories de Picard, n-catégories et autre faune de ce
genre, et ton espoir que je te prouverai peut-être que ces animaux sont tout à
fait indispensables pour faire des maths sérieuses dans telle ou telle
circonstance. J'ai bien \add{peur} que cet espoir ne soit déçu, je crois que
jusqu'à maintenant on a toujours pu s'en tirer en éludant de tels objets et
l'engrenage dans lequel ils pourraient nous entrainer. Est-ce nécessairement une
raison pour continuer à les éluder ? Les situations où on a l'impression
"d'éluder" en effet me semblent en tous cas devenir toujours plus nombreuses —
et si on s'abstenait de tirer une situation complexe et chargée de mystère au
clair, chaque fois qu'on ne serait pas \emph{forcé} de le faire pour des raisons
techniques provenant de la math déjà faite, — il y aurait sans doute beaucoup de
parties des maths aujourd'hui reputées "sérieuses" qui n'auraient jamais
\add{été} développées (Il n'est pas dit non plus que le monde s'en trouverait
plus mal …). Ton commentaire (que j'ai également entendu chez Deligne) que la
classification d'objets géométriques relativement merdiques se réduit finalement
à des invariants cohomologiques essentiellement "bien connus" et relativement
simples n'est pas non plus convainquant ; n'est ce pas négliger la différence
entre la \emph{compréhension d'un objet géométrique}, et la détermination de sa
"classe à isomorphisme (ou équivalence …) près" ?

\page{39}
\note{p. de l'auteur 24.}

Tu me demandes des exemples "convainquants" de 2-catégories de Picard. Voici
quelques \struck{cas} \add{exemples,} en vrac \add{(je ne sais s'ils seront
convainquants !)} :

\marginal{(22)}
1) Si $L$ est un lien \add{(22)} de centre $Z$ sur le topos $X$, les
\emph{gerbes liées par $Z$} forment une 2-catégorie de Picard stricte,
représentée par le complexe $R\Gamma_X(Z)$ tronqué en degré 2, dont les objets
de cohomologie non triviaux sont les $H^i(X,Z)$, $0\leqslant i\leqslant 2$. Les
gerbes liées par $L$ forment un \emph{pseudo-2-torseur} sous la gerbe
précédente, qui est un 2-torseur (i.e.\ non vide) sss une certaine obstruction
dans $H^3(X,Z)$ est nulle. Pour comprendre cette classe de notre point de vue,
il y a lieu de passer aux 2-champs correspondants : le 2-champ de Picard
\add{strict} des $Z$-gerbes \struck{liés par \ill{} et le 2-champ des} sur des
objets variables de $X$, et le 2-champ des $L$-gerbes sur des objets variables.
Ce dernier est bel et bien un 2-torseur sous le champ précédent, or la
classification de ces \struck{deux} 2-torseurs (à 2-équivalence près) se fait
par le $H^3(X,Z)$, \struck{De façon précise} (tout comme les $Z$-\add{1-}gerbes
peuvent être interpétées\note{sic} comme des \uncertain{1-}torseurs sous la
$Z$-1-champ \add{de Picard strict} des $Z$-torseurs, et sont \struck{par
conséquent} classifiées par le $H^2(X,Z)$). On voit déjà, bien sûr, poindre
ici l'oreille de la 3-catégorie de Picard stricte des 2-gerbes liées par $Z$,
ou (de façon équivalente) des 2-torseurs sous le 2-champ de Picard strict des
$Z$-1-gerbes ; cette 3-catégorie de Picard stricte étant décrite par
$R\Gamma_X(Z)$ tronqué en dimension 3, ayant comme invariants de cohomologie
non triviaux les $H^i(X,Z)$ $(0\leqslant i\leqslant 3)$. Quant au 3-champ de
Picard correspondant, il est décrit par une résolution injective de $Z$
tronquée en degré 3, alors que le 2-champ de Picard précédent se décrivait en
tronquant en degré 2.

2. Si $M$ et $N$ sont deux faisceaux abéliens sur $X$, les \emph{champs de
Picard} (NB : 1-champs !) \emph{d'invariants} $M$ \emph{et} $N$ forment
eux-même une \add{2-}catégorie de Picard stricte, représentée sans doute par le
complexe $\mathrm{RHom}(X(M),N)$ \add{(23)} tronqué en \struck{dimension} degré
2, dont les invariants de chomologie\note{sic} non triviaux sont donc "le drôle
de $\mathrm{Ext}^2$" de ma lettre à Deligne, et les Honnètes
$\mathrm{Ext}^i(M,N)$ $(0\leqslant i\leqslant 1)$.
\note{l'exposant de ce premier $\mathrm{Ext}$ ressemble à un 1 ; l'encadrement $0\leqslant i\leqslant 1$ demande $i$.}
Les champs de Picard stricts forment une sous-2-catégorie de Picard pleine,
représentée par $\mathrm{RHom}(M,N)$ tronqué en degré 2, d'invariants les
$\mathrm{Ext}^i(M;N)$ $(0\leqslant i\leqslant 2)$. Bien sûr, $\mathrm{Ext}^2$
donne les \add{0-}objets à équivalence près, $\mathrm{Ext}^1$ les
automorphismes à isomorphisme près de l'objet nul, $\mathrm{Ext}^0$ les
automorphismes de l'automorphisme identique dudit … Je n'ai pas réfléchi à une
bonne interprétation géométrique de la $n$-catégorie de Picard associée à
$\mathrm{RHom}(M,N)$ tronqué en degré $n$, et encore moins bien sûr pour
$\mathrm{RHom}(X(M),N)$, mais sans doute il faut regarder dans la direction des
$n$-champs de Picard.

3. Soit $G$ un Groupe sur $X$, opérant sur un faisceau abélien $N$. Les
\emph{champs en (Gr)-catégories sur $X$ liés par} $(G,N)$ forment une
2-catégorie de Picard,

\page{40}
\note{p. de l'auteur 25.}

dont les invariants sont $H^3(B_G \bmod X,N)$, $H^2(B_G \bmod X,N)$, et
$Z^1(G,N)$ (groupe des 1-cocyles de $G$ à coéfficients dans $N$) — je te laisse
le soin de deviner quel est le complexe tronqué qui le décrit ! J'ai écrit il y
a quelques mois à Deligne à ce sujet, et l'ai prié de t'envoyer une copie de
\struck{\ill{}} la lettre.

4) Soit $X$ un topos localement annelé, on peut considérer les \emph{Algèbres
d'Azumaya sur} $X$ (i.e.\ les Algèbres localement isomorphes à une alg.\ de
matrices d'ordre $n$, $n\geqslant 1$) comme les objets d'une 2-catégorie de
Picard, où la catégorie $\underline{\mathrm{Hom}}(A,B)$, pour $A$ et $B$ des
Algèbres d'Azumaya, est la catégorie des "trivialisations" de $A^{\circ}\otimes
B$, i.e.\ des couples $(E,\varphi)$, $E$ un Module localement libre et
$\varphi$ un isomorphisme $\underline{\mathrm{End}}(E)\simeq A^{\circ}\otimes
B$. Il faut travailler un peu pour définir les accouplements
$\underline{\mathrm{Hom}}(A,B)\times\underline{\mathrm{Hom}}(B,C)\to
\underline{\mathrm{Hom}}(A,C)$ ; l'opération $\otimes$ dans la 2-catégorie de
Picard à construire est bien sûr le produit tensoriel d'Algèbres, et
l'opération "puissance $-1$" est le passage à l'algèbre opposée. On vérifie
\struck{, passant aux champs correspondants, qu'en} qu'en associant à toute
Algèbre d'Azumaya la 1-gerbe de ses trivialisations, on trouve un
2-$\otimes$-foncteur de la 2-catégorie de Picard (dite "de Brauer") dans celle
des 1-gerbes liées par $\underline{G}_m$, qui est 2-fidèle. Les invariants de
la première sont donc les groupes $H^2(X,\underline{G}_m)_{\mathrm{Br}}$,
$H^1(X,\underline{G}_m)$ et $H^0(X,\underline{G}_m)$, où dans le premier terme
l'indice Br désigne le sous-groupe du $H^2$ formé des classes de cohomologie
provenant d'Algèbres d'Azumaya. On aurait envie de parler du 2-champ de Picard
des Algèbres d'Azumaya sur des objets variables de $X$, mais c'est bien une
2-catégorie de Picard fibrée sur $X$, mais pas tout à fait un 2-champ (whatever
that means), sans doute — la condition de 2-recollement (whatever that means)
ne doit pas être satisfaite — sinon il n'y aurait pas d'indice Br au $H^2$ …

\marginal{espace}
(\add{App. 2}) La considération des $n$-catégories de Picard \add{strictes}
(qui s'imposent à nous pas à pas dans un contexte essentiellement
"commutatif") me semblent la clef du passage de l'algèbre homologique ordinaire
("commutative"), en termes de complexes, à une algèbre homologique non
commutative, du fait qu'elles donnent une interprétation géométrique correcte
des "complexes tronqués à l'ordre $n$" \add{(en tant qu'objets de catégories
dérivées)}, donc, essentiellement (par passage à la limite sur $n$) des
complexes tout courts. L'idée naive qui se présente est alors que les
"complexes non commutatifs" (qui seraient les objets-fantômes d'une algèbre
homologique non commutative) sont peut-être ce qui reste des $n$-catégories de
Picard (strictes) quand on oublie leur caractère additif, i.e.\ leur structure
de Picard — c'est à dire qu'on ne retient que la $n$-catégorie !
\add{(Quand on se place sur un topos $X$, on s'intéresse donc aux $n$-champs
sur $X$ …)} A vrai dire, cette idée est venue d'abord d'une autre direction,
quand il s'est agi en géométrie algébrique de démontrer des \struck{formules}
\add{théorèmes} de Lefschetz

\note{La phrase et la lettre se poursuivent à la page suivante (lot 03).}

\end{document}
